\documentclass[8pt,mathserif,professionalfont]{beamer}
\usepackage[utf8]{inputenc}   % pacote para acentuao
\usepackage[scaled]{helvet}
\usepackage{tikz}
\usepackage{varwidth} 

\definecolor{LightButter}{rgb}{0.98,0.91,0.31}
\definecolor{LightOrange}{rgb}{0.98,0.68,0.24}
\definecolor{LightChocolate}{rgb}{0.91,0.72,0.43}
\definecolor{LightChameleon}{rgb}{0.54,0.88,0.20}
\definecolor{LightSkyBlue}{rgb}{0.45,0.62,0.81}
\definecolor{LightPlum}{rgb}{0.68,0.50,0.66}
\definecolor{LightScarletRed}{rgb}{0.93,0.16,0.16}
\definecolor{Butter}{rgb}{0.93,0.86,0.25}
\definecolor{Orange}{rgb}{0.96,0.47,0.00}
\definecolor{Chocolate}{rgb}{0.75,0.49,0.07}
\definecolor{Chameleon}{rgb}{0.45,0.82,0.09}
\definecolor{SkyBlue}{rgb}{0.20,0.39,0.64}
\definecolor{Plum}{rgb}{0.46,0.31,0.48}
\definecolor{ScarletRed}{rgb}{0.80,0.00,0.00}
\definecolor{DarkButter}{rgb}{0.77,0.62,0.00}
\definecolor{DarkOrange}{rgb}{0.80,0.36,0.00}
\definecolor{DarkChocolate}{rgb}{0.56,0.35,0.01}
\definecolor{DarkChameleon}{rgb}{0.30,0.60,0.02}
\definecolor{DarkSkyBlue}{rgb}{0.12,0.29,0.53}
\definecolor{DarkPlum}{rgb}{0.36,0.21,0.40}
\definecolor{DarkScarletRed}{rgb}{0.64,0.00,0.00}
\definecolor{Aluminium1}{rgb}{0.93,0.93,0.92}
\definecolor{Aluminium2}{rgb}{0.82,0.84,0.81}
\definecolor{Aluminium3}{rgb}{0.73,0.74,0.71}
\definecolor{Aluminium4}{rgb}{0.53,0.54,0.52}
\definecolor{Aluminium5}{rgb}{0.33,0.34,0.32}
\definecolor{Aluminium6}{rgb}{0.18,0.20,0.21}
\definecolor{Black}{rgb}{0,0,0}

\tikzstyle{box1} = [
   draw=ScarletRed,
   very thick,
   rectangle, 
   inner sep=10pt, 
   inner ysep=4pt
]

\tikzstyle{box2} = [
   draw=ScarletRed,
   fill=ScarletRed,
   opacity=0.4,
   rectangle, 
   inner sep=10pt, 
   inner ysep=4pt
]

\tikzstyle{title1} = [
   rectangle, 
   very thick,
   fill=ScarletRed, 
   draw=ScarletRed,
   text=white,
   font=\bf, 
   %right=0.5cm, 
   below right,
   execute at begin node={\begin{varwidth}{8.5cm}},
   execute at end node={\end{varwidth}}
]

\tikzstyle{title2} = [
   rectangle, 
   very thick,
   fill=white,
   draw=ScarletRed,
   text=black,
   font=\bf, 
   %right=0.5cm, 
   below right,
   execute at begin node={\begin{varwidth}{8.5cm}},
   execute at end node={\end{varwidth}}
]

\tikzstyle{op9} = [draw=white, fill=white, opacity=0.9]
\tikzstyle{op8} = [draw=white, fill=white, opacity=0.8]
\tikzstyle{op7} = [draw=white, fill=white, opacity=0.7]
\tikzstyle{op6} = [draw=white, fill=white, opacity=0.6]
\tikzstyle{op5} = [draw=white, fill=white, opacity=0.5]
\tikzstyle{op4} = [draw=white, fill=white, opacity=0.4]
\tikzstyle{op3} = [draw=white, fill=white, opacity=0.3]
\tikzstyle{op2} = [draw=white, fill=white, opacity=0.2]
\tikzstyle{op1} = [draw=white, fill=white, opacity=0.1]

\usetheme{Antibes}
%\usetheme{Dresden}
%\usetheme{Frankfurt}
%\usetheme{Bergen}
%\useoutertheme{shadow}

\usecolortheme[RGB={170,0,0}]{structure}

\setbeamertemplate{blocks}[rounded][shadow=true]
\setbeamertemplate{title page}[default][colsep=-4bp,rounded=true,shadow=true]
\beamertemplatesquareitem
\setbeamercovered{dynamic}
%\usefonttheme{serif}
%\usefonttheme{structurebold}
\usepackage{natbib}
\usepackage{graphicx}
\usepackage{mathptmx}      % use Times fonts if available on your TeX system
\usepackage{amsmath}
\usepackage{amsfonts}
\usepackage{multicol}
\usepackage{multirow}
\usepackage{booktabs}
\usepackage{float}
\usepackage{natbib}
\usepackage{mdwlist}
\usepackage{booktabs}
\usepackage{url}
\usepackage{rotating}
\usepackage{tabularx}
\usepackage[symbol*]{footmisc}
\usepackage{bibentry}

\setbeamertemplate{navigation symbols}{}
\setbeamertemplate{footline}[frame number]

%\renewcommand\mathfamilydefault{\rmdefault}

\newcommand{\mc}[3]{
  \multicolumn{#1}{#2}{#3}
}

\newcommand{\mr}[2]{
  \multirow{#1}{*}{#2}
}
\usepackage{algorithmic}

\def\Beginboxit
   {\par
    \vbox\bgroup
       \hrule
       \hbox\bgroup
          \vrule \kern1.2pt 
          \vbox\bgroup\kern1.2pt
   }

\def\Endboxit{
                  \kern1.2pt
               \egroup
          \kern1.2pt\vrule
        \egroup
       \hrule
     \egroup
   }
\newenvironment{boxit}{\Beginboxit}{\Endboxit}
\newenvironment{boxit*}{\Beginboxit\hbox to\hsize{}}{\Endboxit}
\newcommand{\labelcaption}[2]{
    \caption{#2}
    \label{#1}
}
\newcommand{\iniciaproc}[2]{
    \smallskip
    \textbf{Algoritmo}~\textit{#1}~(#2)~
    \begin{algorithmic}[1]
}

%label, caption
\newcommand{\fimproc}[2]{
    \end{algorithmic}
}



\begin{document}
\title{Metaheurística \textit{Late Acceptance Hill Climbing} aplicada ao Problema de Escalonamento de Enfermagem}
\author{\textbf{Victória S. Portella\inst{1}}\and {Toni I. Wickert\inst{1}}\and {Luciana S. Buriol\inst{1}}}
\titlegraphic{\includegraphics[scale=0.4]{inf-20}}
\institute[Instituto de Informática - UFRGS]{
   \inst{1}
   Instituto de Informática \\ Universidade Federal do Rio Grande do Sul \\
}
\date[]{}
\subject{}



\AtBeginSection[]%
{%
   \begin{frame}%[plain] 
        \begin{center}%
        {
           \usebeamercolor[bg]{frametitle}
           \usebeamerfont{section title}
           \Huge
           \textbf{\insertsection}
        }
        \insertsubsectionnavigation{10cm}
        \end{center}%
     \end{frame}
}


%-------------------------------------------------------------------------------------------------
\frame{
   \titlepage
}

%-------------------------------------------------------------------------------------------------
\frame{
   \frametitle{Roteiro}
   \setcounter{tocdepth}{1}

   %XXX Ativar na versão final \tableofcontents[pausesections]
   \tableofcontents
}
%-------------------------------------------------------------------------------------------------
\section{Introdução}

\begin{frame}{Introdução} % o qq é o problema para leigos, pq é um problema, pra q ->motivação, onde ocorre,
									% com quem ocorre. qual objetivo do problema minimizar problema.

   O processo de elaboração de \emph{escalas de enfermagem} é um problema operacional geralmente encontrado no ambiente hospotalar. 
   \medskip
   
   O objetivo desse processo é definir os turnos de trabalho de cada enfermeiro considerando um horizonte de planejamento e um determinado conjunto de restrições:
   	
   \begin{itemize}
      \item Restrições impostas por lei.
      \item Normas particulares da instituição.
      \item Preferências dos funcionários.
   \end{itemize}
   \medskip
   Uma escala deficiente pode trazer consequências negativas para os pacientes e funcionários.


\end{frame}

%--------------------------------------------------------------------------------------------------
\begin{frame}{Exemplo de escala de enfermagem} % escala completa
   \begin{center}
   \includegraphics[scale=0.55]{tabelaGeral.pdf} 
   \end{center}
\end{frame}

%--------------------------------------------------------------------------------------------------
\begin{frame}{Trabalhos relacionados}
	
		%\item Nas últimas décadas vários métodos foram propostos para resolver diversas variantes do problema, mas é difícil compará-los.
	\textbf{\textit{International Nurse Rostering Competition} (INRC-I)} \citep{haspeslagh2010first}.
		
		\begin{itemize}
			\item \structure{Problema de um único estágio}
			\item Programação Inteira~\citep{valouxis2012systematic}
			\item Programação por Restrições~\citep{nonobe2010}
			\item \textit{Branch and Price e Ejection Chain}~\citep{burke10}
		\end{itemize}
		
		\bigskip
		
	\textbf{\textit{Second International Nurse Rostering Competition} (INRC-II)} \citep{ceschia2015model}
		
		\begin{itemize}
			\item \structure{Problema multi-estágio}
			\item Programação Inteira Mista baseada em fluxo de rede \textit{multi-commodity} \citep{romer2015multi};
			\item Hiper-heurística \citep{kheirisequence}
			\item Geração de Colunas combinada com VNS \citep{gomes2017variable}
			
		\end{itemize}
	
\end{frame}

%--------------------------------------------------------------------------------------------------
\section{Definição do Problema}

%--------------------------------------------------------------------------------------------------
\begin{frame}[plain]{Problema multi-estágio} 

\begin{center}
   \includegraphics[scale=0.7]{escalageral.pdf} 
   \end{center}
   
\begin{tikzpicture}[overlay]
	\draw<2>[box2] (6.53,3.47) rectangle (8.04,3.88) ;
	\draw<2>[box2] (2.76,1.59) rectangle (4.26,2) ;
	
	\draw<3>[box2] (4.3,6.35) rectangle (7.3,6.75) ;
	\draw<3>[box2] (5.0,4.5) rectangle (8.,4.9) ;
    \draw<3>[box2] (5.0,2.6) rectangle (8.,3) ;
    \draw<3>[box2] (3.5,0.75) rectangle (2.8,1.15) ;
    \draw<3>[box2] (7.25,0.75) rectangle (4.25,1.15) ;

   
\end{tikzpicture}   
   
\end{frame}


%-------------------------------------------------------------------------------------------------
\begin{frame}{Definição das restrições da INRC-II}
\smallskip

\textbf{Restrições Fortes}
\begin{itemize}
\item[\textbf{H1}] O enfermeiro deve trabalhar apenas um turno por dia.
\item[\textbf{H2}] A cobertura mínima deve ser respeitada.
\item[\textbf{H3}] A sucessão de turnos deve ser válida.
\item[\textbf{H4}] Qualificações atribuídas devem pertencem ao conjunto de qualificações do enfermeiro.
\end{itemize}
\medskip
\textbf{Restrições Fracas}
\begin{itemize}
\item[\textbf{S1}] A cobertura ótima deve ser atendida tanto quanto possível.
\item[\textbf{S2ab}] Número mínimo/máximo de dias consecutivos trabalhados.
\item[\textbf{S2cd}] Número mínimo/máximo de alocações consecutivas no mesmo turno.
\item[\textbf{S3ab}] Número mínimo/máximo de dias consecutivos de folga.
\item[\textbf{S4ab}] Enfermeiro prefere não trabalhar em um determinado dia/turno.
\item[\textbf{S5}] Preferência por fins de semana completo.
\item[\textbf{S6ab}] Número mínimo/máximo de dias trabalhados.
\item[\textbf{S7}] Número máximo de fins de semana trabalhados.

\end{itemize}

\begin{tikzpicture}[overlay]
	\draw<2>[box1] (-0.3,1.2) rectangle (10,7) 
       node[title1] at (-0.3,1.2) {Restrições locais: envolvem dados do estágio atual + histórico};
    \draw<3>[box1] (-0.3,0.3) rectangle (10,1.2) 
       node[title1] at (-0.3,1.8) {Restrições globais: envolvem dados de todos os estágios};
\end{tikzpicture}

\end{frame}

%---------------------------------------------------------------------------------------------------
\begin{frame}{Definição formal das restrições utilizando formulação matemática}

Proposta de um modelo não linear mais fácil de compreender que o modelo definido na competição que oferece dois pontos de vista do problema:
\medskip

	\begin{itemize}
		\item Modelo $\mu_1$: 
			\begin{itemize}
				\item Visão de um \structure{único estágio}.
				\item Compreende apenas \structure{restrições locais}.
			\end{itemize}
		
		\smallskip
		\item Modelo $\mu_2$:
		\begin{itemize}
				\item Visão de \structure{vários estágios}.
				\item Expande o modelo $\mu_1$ adicionando \structure{restrições globais}.
			\end{itemize}
	\end{itemize}


\end{frame}

%---------------------------------------------------------------------------------------------------
\begin{frame}{Modelo $\mu_1$: apenas um estágio (Parte I)}

\begin{samepage}
\begin{align}
   \textbf{min~~} 
   Z_e =~~&W^{S1} \sum_{d \in D}\sum_{s \in S}\sum_{k \in K} C^{S1}_{dsk}    \notag +
          W^{S2ab}  \sum_{n \in N}\sum_{(i,j) \in B} (C^{S2a}_{nij} + C^{S2b}_{nij}) + \\ \notag
  		& W^{S2cd}  \sum_{n \in N}\sum_{(i,j) \in B}\sum_{s \in S'} (C^{S2c}_{nijs} + C^{S2d}_{nijs}) +
		  W^{S3}    \sum_{n \in N}\sum_{(i,j) \in B} (C^{S3a}_{nij} + C^{S3b}_{nij}) + \\ \notag
		& W^{S4}    \sum_{n \in N}\sum_{d \in D}\sum_{s \in S'}\sum_{k \in K} (x^e_{ndsk} * P_{nds} + x^e_{ndsk} * P_{nd}) + 
		W^{S5}    \sum_{n \in N} V_{n} * ( 1- \prod_{d \in W} \sum_{s \in S'}\sum_{k \in K} x^e_{ndsk}) -\delta^e
\end{align}  

\begin{align}
		& \sum_{s \in S} \sum_{k \in K} x^e_{ndsk} = 1    && \forall n \in N, d \in D \label{eH1} \\
& \sum_{n \in N} x^e_{ndsk} \geq V^-_{dsk}   && \forall d \in D, s \in S, k \in K \label{eH2} \\
& \sum_{k \in K} (x^e_{n,d-1,s_1,k} + x^e_{n,d,s_2,k}) \leq 1 && _{\forall n \in N, d \in D\setminus \{1\}, (s_1,s_2) \in F} \label{eH3a} \\
& \sum_{k \in K} (x^{e-1}_{n,d,s_1,k} + x^e_{n,1,s_2,k}) \leq 1 && _{\forall n \in N, d = |D|, (s_1,s_2) \in F}\\
& x^e_{ndsk} = 0   && _{\forall n \in N, d \in D, s \in S, k \in K\setminus K_n} \label{eH4} \\ \notag
\end{align}  
\end{samepage}

\begin{tikzpicture}[overlay]
	\draw<2>[box1] (-0.1,4.6) rectangle (11.5,7.3) 
       node[title1] at (-0.1,4.6) {Função objetivo penalizando as restrições S1 a S5.};
    \draw<3>[box1] (1.2,3.8) rectangle (11.1,4.6) 
       node[title1] at (1.2,3.8) {Restrição H1};
    \draw<4>[box1] (1.2,3.0) rectangle (11.1,3.8) 
       node[title1] at (1.2,3.0) {Restrição H2};
    \draw<5>[box1] (1.2,1.6) rectangle (11.1,3.1) 
       node[title1] at (1.2,1.6) {Restrição H3};
    \draw<6>[box1] (1.2,1.0) rectangle (11.1,1.6) 
       node[title1] at (1.2,1) {Restrição H4};
\end{tikzpicture}

\end{frame}
%-------------------------------------------------------------------------------------------
%TODO nao está entrando...
\begin{frame}{Modelo $\mu_1$: apenas um estágio (Parte II)}
\begin{samepage}
\begin{align}
& C^{S1}_{dsk} = \max\{ V^*_{dsk} - \sum_{n \in N} x^e_{ndsk}, 0\} && \forall d \in D, s \in S', k \in K \label{eS1} \\
& b^e_{nuv} \texttt{=} \max \{0, \hspace{-1mm}\sum_{u \leq d \leq v} \sum_{s \in S'} \sum_{k \in K} x^e_{ndsk} \texttt{-} (v \texttt{-}u) \texttt{-}\hspace{-1mm} \sum_{(i,j) \in B_{uv}}\hspace{-1mm} b^e_{nij}\} && \forall n \in N, (u,v) \in B \label{eS2aux1} \\
& b^e_{nuvs} \texttt{=} \max \{0, \sum_{u \leq d \leq v} \sum_{k \in K} x^e_{ndsk} \texttt{-} (v \texttt{-}u) \texttt{-} \hspace{-1mm}\sum_{(i,j) \in B_{uv}} \hspace{-1mm} b^e_{nijs}\} && \forall n \in N, (u,v) \in B, s \in S \label{eS2aux2} \\
& ^eC^{S2a}_{nij} \texttt{=}  \max\{0, b^e_{nij} * ((j - i +1 + h^e_{n}*\lfloor 1/i \rfloor) - L^+_{n})\} && \forall n \in N, (i,j) \in B \label{eS2a} \\
& ^eC^{S2b}_{nij} \texttt{=}  \max\{0, b^e_{nij} * (L^-_{n} - (j - i +1  + h^e_{n}*\lfloor 1/i \rfloor))\} && \forall n \in N, (i,j) \in B \label{eS2b} \\
& ^eC^{S2c}_{nijs} \texttt{=}  \max\{0, b^e_{nijs} * ((j - i +1  + h^e_{ns}*\lfloor 1/i \rfloor) - L^+_{ns})\} && \forall n \in N, (i,j) \in B, s \in S'\label{eS2c} \\
& ^eC^{S2d}_{nijs} \texttt{=}  \max\{0, b^e_{nijs} * (L^-_{ns} - (j - i +1 + h^e_{ns}*\lfloor 1/i \rfloor))\} && \forall n \in N, (i,j) \in B, s \in S' \label{eS2d} \\
& ^eC^{S3a}_{nij} \texttt{=}  \max\{0, b^e_{nijs} * ((j - i +1  + h^e_{ns}*\lfloor 1/i \rfloor) - G^+_{n})\} && \forall n \in N, (i,j) \in B, s \texttt{=} |S|\label{eS3a} \\
& ^eC^{S3b}_{nij} \texttt{=}  \max\{0, b^e_{nijs} * (G^-_{n} - (j - i +1  + h^e_{ns}*\lfloor 1/i \rfloor))\} && \forall n \in N, (i,j) \in B, s \texttt{=} |S| \label{eS3b} \\
& x^e_{ndsk} \in \{0,1\}  &&  \forall n \in N, d \in D, s \in S, k \in K  \notag  \\
& b^e_{nijs},  b^e_{nij} \in \{0,1\}  &&  \forall n \in N, (i,j) \in B, s \in S \notag \\
& ^eC^{S2a}_{nij}, ^eC^{S2b}_{nij}, ^eC^{S2c}_{nijs}, ^eC^{S2d}_{nijs}, ^eC^{S3a}_{nij}, ^eC^{S3b}_{nij} \in \mathbb{N^*} && \forall n \in N, (i,j) \in B, s \in S \notag 
\end{align}
\end{samepage}

\begin{tikzpicture}[overlay]
	\draw<2>[box1] (0.4,6.7) rectangle (11.0,7.4) 
       node[title1] at (0.4,6.7) {Restrição S1};
    \draw<3>[box1] (0.4,6.7) rectangle (11.0,5.2) 
       node[title1] at (0.4,5.2) {Restrições auxiliares adicionadas no modelo para identificar \\ blocos de alocações consecutivas};
    \draw<4>[box1] (0.4,5.2) rectangle (11.0,4.1) 
       node[title1] at (0.4,4.1) {Restrições S2a e S2b};
    \draw<5>[box1] (0.4,4.1) rectangle (11.0,3.1) 
       node[title1] at (0.4,3.1) {Restrições S2c e S2d};
    \draw<6>[box1] (0.4,3.1) rectangle (11.0,2.1) 
       node[title1] at (0.4,2.1) {Restrições S3a e S3b};
   % \draw<7>[box1] (0.4,2.1) rectangle (11.0,1.6) 
    %   node[title1] at (0.4,1.6) {Variáveis de decisão};
    %\draw<8>[box1] (0.4,1.7) rectangle (11.0,0.6) 
     %  node[title1] at (0.4,2.2) {Variáveis auxiliares};
    
\end{tikzpicture}

\end{frame}

%---------------------------------------------------------------------------------

\begin{frame}{Modelo $\mu_2$: considera todos os estágios ao mesmo tempo}

\begin{align}
   \textbf{min~~}  &Y =  \sum_{e \in E} Z_e + \sum_{n \in N} (W^{S6} C^{S6a}_n + W^{S6} C^{S6b}_n + W^{S7} C^{S7}_n) 
\end{align}                     
\begin{align}
& C^{S6a}_n =  \max\{0, \sum_{d \in D} \sum_{s \in S'}\sum_{k \in K} x^e_{ndsk} - Q^+_{n}\} && \forall e \in E, n \in N \label{eS6a} \\
& C^{S6b}_n =  \max\{0, Q^-_{n} - \sum_{d \in D} \sum_{s \in S'}\sum_{k \in K} x^e_{ndsk}\}  && \forall e \in E, n \in N \label{eS6b} \\
& C^{S7}_n =  \max\{0,\sum_{e \in E}\max_{d \in W, k \in K, s \in S'} (x^e_{ndsk}) - R^+_{n}\}  && \forall n \in N \label{eS7} \\
& C^{S6a}_{n}, C^{S6b}_{n}, C^{S7}_{n} \in \mathbb{N} && \forall n \in N \notag
\end{align}

\begin{tikzpicture}[overlay]
	\draw<2>[box1] (1.3,3.6) rectangle (11.0,4.5) 
       node[title1] at (1.3,3.6) {Função objetivo};
    \draw<3>[box1] (1.3,1.6) rectangle (11.0,3.3) 
       node[title1] at (1.3,1.6) {Restrição S6};
    \draw<4>[box1] (1.3,1.0) rectangle (11.0,1.7) 
       node[title1] at (1.3,1.0) {Restrição S7};
    %\draw<5>[box1] (1.3,0.4) rectangle (11.0,1.0) 
     %  node[title1] at (1.3,0.4) {Variáveis auxiliares};
   
\end{tikzpicture}
\end{frame}

%gancho: problema q simule o m2 considerando informações de apenas um estágio, especificamente cabe estimar o calculo das restricoes s7 e s8

%---------------------------------------------------------------------------------------------------
\section{Algoritmo Proposto}


%---------------------------------------------------------------------------------------------------
\begin{frame}{\textit{Late Acceptance Hill Climbing} (\cite{burke2012late})}
%bloco 1: solução incial : gerada aleatoriamente, garantindo h1 e h4
%bloco 2: geraçao da sol candidata
%bloco 3: função objetivo

  \iniciaproc{LAHC}{$t$}
	  \STATE Gera solução inicial $s$
      \STATE $s^* \gets s$ 
      \FOR{$w \in 0...t-1$}
      	\STATE $L_w \gets C(s)$
      \ENDFOR
      \STATE $i \gets 0$ 
      \REPEAT 
         \STATE Gera uma solução candidata $s'$
         \STATE $v \gets i$ $mod$ $t$
         \IF {$C(s') \leq L_v$ ou $C(s') \leq C(s)$}
            \STATE $s \gets s'$
            \IF {$C(s) \le C(s^*)$}
           		\STATE $s^* \gets s$
         	\ENDIF
         \ENDIF
         \STATE $L_v \gets C(s)$
         \STATE $i \gets i+1$
      \UNTIL{condição de parada seja satisfeita}
      \RETURN $s^*$
   \fimproc{alg:tabuSearch}{Pseudo-código do LAHC}

\begin{tikzpicture}[overlay]
	\draw<2>[box1] (0.4,6.3) rectangle (3.3,6.7) 
       node[title1] at (0.4,6.3) {Geração da Solução Inicial};
    \draw<3>[box1] (0.5,4.35) rectangle (4.6,3.9) 
       node[title1] at (0.5,3.9) {Geração da Solução Candidata};
    \draw<4>[box1] (1.4,5.7) rectangle (2.1,5.25) 
       node[title1] at (1.4,5.25) {Função objetivo custo C(.)};
   
\end{tikzpicture}	
   

\end{frame}


%-----------------------------------------------------------------------------------------------
\begin{frame}{Geração da solução candidata}
	Solução $s'$ é escolhida a partir de uma estrutura de vizinhança $\mathcal{N}(s)$ onde $s$ é a solução atual.
	A estrutura $\mathcal{N}(s)$ possui 7 movimentos que compõe a vizinhança, de $M_1$ a $M_7$:
	\medskip
	\begin{enumerate}
		\item Movimento 1 ($M_1$): Modificação de um turno;
		\item Movimento 2 ($M_2$): Troca de turnos em blocos;
		\item Movimento 3 ($M_3$): Modificação de qualificação;
		\item Movimento 4 ($M_4$): Atribuição de uma folga;
		\item Movimento 5 ($M_5$): Fim de semana completo;
		\item Movimento 6 ($M_6$): Troca de turno de trabalho com turno de folga;
		\item Movimento 7 ($M_7$): Troca de blocos para um enfermeiro.
	\end{enumerate}

	\medskip
	
	\structure{Seleção do movimento:} Selecionados de acordo com uma dada probabilidade.
	

\end{frame}

%-----------------------------------------------------------------------------------------------
\begin{frame}{Função objetivo do LAHC usada em um estágio $e$}

\begin{align}
& C(.) = Z_e + \sum_{n \in N} \left(C^{S6a}(n,e) +  C^{S6b}(n,e) +  C^{S7}(n,e) \right)   
\end{align}

\begin{align}
& Q^{+}(n,e) = \max\{0, \lceil (( Q^+_{n} - \beta(n,e)) / (|E|-e) + 0.5)\rceil\} \label{ests6a}\\
& Q^{-}(n,e) = \lfloor ( Q^-_{n} - \beta(n,e)) / (|E|-e) \rfloor   \label{ests6b} \\
&R^{+}(n,e) = \max\{0, \lceil (( R^+_{n} - \gamma(n,e)) / (|E|-e) +0.5)\rceil\}  \label{rs7}
\end{align}


\begin{align}
& C^{S6a}(n,e) = max\{0, \sum_{d \in D} \sum_{s \in S'}\sum_{k \in K} x^e_{ndsk} -  Q^{+}(n,e)\}  \label{cs6a}\\
& C^{S6b}(n,e) = max\{0,  Q^{-}(n,e) - \sum_{d \in D} \sum_{s \in S'}\sum_{k \in K} x^e_{ndsk}\}  \label{cs6b} \\
& C^{S7}(n,e) = max\{0, \sum_{d \in D} \sum_{s \in S'}\sum_{k \in K} x^e_{ndsk} -  R^{+}(n,e)\} 
\end{align}

\begin{tikzpicture}[overlay]
	\draw<1>[box1] (2,6.2) rectangle (8.7,5.1)
	 node[title1] at (2,5.1) {Função objetivo};      
	\draw<2>[box1] (2,4.8) rectangle (8.7,3.1)
	 node[title1] at (2,3.1) {Novos limites locais para restrições globais};
	 \draw<3>[box1] (2.2,2.8) rectangle (8.7,0.1)
	 node[title1] at (2.2,0.1) {Estimativa das restrições S6 e S7};        
	\end{tikzpicture}	

\end{frame}

%--------------------------------------------------------------------------------------------------
\begin{frame}[plain]{Exemplo de estimativa da restrição S6a considerando $Q^+_{3} = 16$} 


	\hspace{-0.4cm}
   \includegraphics[scale=0.7]{estimativa2.pdf} 

   
\begin{tikzpicture}[overlay]	
	\draw<1>[box2] (4.6,6.1) rectangle (6.1,6.5) ;
	\draw<1>[box2] (3.85,4.2) rectangle (6.85,4.6) ;
    \draw<1>[box2] (3.85,2.3) rectangle (6.85,2.75) ;
    \draw<1>[box2] (1.6,0.45) rectangle (2.35,0.9) ;
    \draw<1>[box2] (3.1,0.45) rectangle (6.1,0.9);
    \draw<1>[box1] node[title2] at (7.1,7.) {$Q^{+}(3,0) = \lceil16/4+0.5\rceil = 4$};
    \draw<1>[box1] node[title2] at (7.1,5.15) {$Q^{+}(3,1) = \lceil14/3+0.5\rceil = 5$};
    \draw<1>[box1] node[title2] at (7.1,3.28) {$Q^{+}(3,2) = \lceil10/2+0.5\rceil = 5$};
    \draw<1>[box1] node[title2] at (7.1,1.4) {$Q^{+}(3,3) = \lceil6/1+0.5\rceil = 6$};

   
\end{tikzpicture}   
   
\end{frame}

%---------------------------------------------------------------------------------------------------
\section{Experimentos computacionais}
%formato de pergunta
\begin{frame}{Questões de pesquisa}
	\begin{enumerate}
		\item Qual a combinação de parâmetros é mais apropriada para o algoritmo proposto?
		\smallskip
		\item Como o algoritmo se compara com os demais competidores?
		\smallskip
		\item O método proposto conseguiria ser classificado para a final?
		\smallskip
		\item O método proposto é robusto quanto à reprodução dos resultados?
		
	\end{enumerate}
\end{frame}

%---------------------------------------------------------------------------------------------------
\begin{frame}{Ambiente experimental}
	\begin{itemize}
		\item \textit{Hardware}: computador servidor com um processador Intel Xeon 2.83GHz x 4, 8GB RAM. Utilização de uma única \textit{thread};
		\item Sistema Operacional Linux Ubuntu 16.04 LTS - 64 bits;
		\item Linguagem de programação C++, compilador g++ (v. 5.4.0);
		\item Computação do processo de variação de amostras foi usado \structure{\textit{Mersenne Twister}};
		\item Resultados do experimentos foram certificados pelo \structure{validador da competição};
		\item 15 competidores internacionais, incluindo \structure{empresas e grupos de pesquisa};
		\item \structure{Tempo limite} para cada semana determinado pelo \textit{benchmark}:
		
	\end{itemize}
	
	\begin{table}[ht!]		
		\label{tab:instances}
		\centering
		\scalebox{0.8}{
		\begin{tabular}{r|rrrrrrrrrr}
		\textit{Enfermeiros} & 30 & 35 & 40 & 50 & 60 & 70 & 80 & 100 & 110 & 120 \\
		\hline
		\textit{Tempo(s)} &  59.7 & 82.0 & 104.5 & 149.3 & 194.1 & 238.5 & 283.7 & 373.3 & 417.4 & 462.0    \\
		\end{tabular}
		}
	\end{table}

\end{frame}
%-------------------------------------------------------------------------------------------
\begin{frame}{Conjunto de instâncias da INRC-II}

	\textbf{Instâncias \textit{Late}:}
	\begin{itemize}
		\item Conjuntos de enfermeiros: 30, 40, 50, 60, 80, 100 e 120;
		\item Horizonte de planejamento: 4 e 8 semanas;
		\item Total: 28 instâncias.
	\end{itemize}
	\medskip
	\textbf{Instâncias \textit{Hidden}:}
	\begin{itemize}
		\item Conjunto de enfermeiros: 35, 70 e 110;
		\item Horizonte de planejamento: 4 e 8 semanas;
		\item Total: 60 instâncias.
	\end{itemize}


	
\end{frame}

%-------------------------------------------------------------------------------------------
\begin{frame}{Resultados do ajuste de parâmetros do LAHC com a ferramenta I-Race}

\textbf{Objetivo:}
\begin{itemize}
	\item Definir parâmetro $TL$ (Tamanho da Lista).
	\item Definir os parâmetros $P_{M1}$~a~$P_{M7}$ (probabilidade dos movimentos).
\end{itemize}

\textbf{Configurações:}
\begin{itemize}
	\item Usamos instâncias \textit{Late} como treino.
	\item \textit{Budget} = 1000.
	
\end{itemize}

\begin{table}[ht!]
\centering
\small
\begin{tabular}{llr}
\toprule
    \textbf{Parâmetros} & \multicolumn{1}{c}{\textbf{Intervalo}} & \multicolumn{1}{c}{\textbf{Melhores parâmetros}}\\
     \hline
    $TL$ & \{5000, 10000, 15000, 20000\} & 5000\\
    $P_{M1}$ & \{0 ... 10\} & 8 (19\%)\\
    $P_{M2}$ & \{0 ... 10\} & 9 (21\%)\\
    $P_{M3}$ & \{0 ... 10\} & 8 (19\%)\\
    $P_{M4}$ & \{0 ... 10\} & 10 (23\%)\\
    $P_{M5}$ & \{0 ... 10\} & 0 (0\%)\\
    $P_{M6}$ & \{0 ... 10\} & 1 (2\%)\\
    $P_{M7}$ & \{0 ... 10\} & 7 (16\%)\\
   \bottomrule
\end{tabular}
\label{tab:iraceParam}
\end{table}
	
	\begin{tikzpicture}[overlay]
	\draw<2>[box1] (1.5,3) rectangle (3.2,0.3); 
 	\draw<3>[box1] (3.2,3) rectangle (6.7,0.3); 
    \draw<4>[box1] (6.7,3) rectangle (9.5,0.3); 
   
	\end{tikzpicture}	
	
\end{frame}


%-----------------------------------------------------------------------------------------------
\begin{frame}{Resultados da etapa de classificação da INRC-II}

\begin{table}[H]
\scalebox{0.52}{
\definecolor{brown}{rgb}{0.28, 0.02, 0.03}
\renewcommand\it{\color{brown}\bf$^*$}
\centering
\label{tab:late}
\begin{tabular}{lrrrrrrrrrrrrrrrr}
  \toprule
id & LAHC & $A_1$ & $A_2$ & $A_3$ & $A_4$ & $A_5$	 & $A_6$ & $A_7$ & $A_8$ & $A_9$ & $A_{10}$ & $A_{11}$ & $A_{12}$ & $A_{13}$ & $A_{14}$ & $A_{15}$\\ 
   \midrule
  1 &  24.07 &   0 & 2.01 & 10.60 & 15.19 & 14.61 & 18.62 & 18.34 & \it 36.68 & \it 34.96 & \it 42.69 & \it 87.97 & \it 54.15 & \it 89.11 & \it 225.50 & \it 464.50 \\
    2 & 21.45 &   0 & 0.78 & 10.08 & 10.34 & 8.53 & 11.63 & 15.25 & \it 24.81 & \it 36.43 & \it 38.50 & \it 69.25 & \it 53.23 & \it 72.09 & \it 229.20 & \it 448.10 \\
    3 & 37.04 &   0 & 1.31 & 28.11 & 28.32 & 13.72 & 34.20 & 36.17 & \it 56.21 & \it 72.77 & \it 75.82 & \it 94.99 & \it 79.74 & \it 115.47 & \it 307.60 & \it 823.10 \\
    4 & 42.63 &   0 & 4.74 & 25.26 & 36.05 & 24.74 & 38.95 & 38.42 & \it 49.47 & \it 81.05 & \it 82.63 & \it 90.26 & \it 112.89 & \it 169.47 & \it 359.20 & \it 1012.90 \\
    5 & 17.39 &   0 & 2.32 & 9.86 & 8.41 & \it 20.29 & \it 22.90 & \it 31.88 & \it 25.22 & \it 48.12 & \it 59.42 & \it 87.25 & \it 71.88 & \it 135.65 & \it 391.00 & \it 751.00 \\
    6 & 18.06 &   0 & 5.24 & 13.35 & 8.90 & 17.02 & \it 22.51 & \it 23.82 & \it 35.08 & \it 41.88 & \it 61.26 & \it 77.23 & \it 72.51 & \it 102.09 & \it 374.60 & \it 657.10 \\
    7 & 49.76 &   0 & 6.76 & 20.93 & 21.26 & 30.43 & \it 49.76 & 45.09 & 46.05 & \it 59.10 & \it 92.11 & \it 94.20 & \it 118.68 & \it 142.35 & \it 421.30 & \it 1027.50 \\
    8 & 35.88 &   0 & -1.18 & 20.60 & 15.28 & 24.09 & \it 42.69 & \it 37.54 & 26.91 & \it 57.97 & \it 78.74 & \it 88.54 & \it 101.33 & \it 147.01 & \it 398.70 & \it 996.30 \\
    9 & 45.90 &   0 & 6.89 & 20.98 & 13.44 & 23.93 & 33.77 & 40.00 & \it 60.98 & \it 73.77 & \it 94.75 & \it 120.33 & \it 111.48 & \it 160.66 & \it 537.40 & \it 1063.60 \\
   10 & 40.54 &   0 & 11.15 & 26.69 & 20.61 & 32.09 & \it 44.59 & \it 47.64 & \it 64.53 & \it 84.80 & \it 106.08 & \it 101.01 & \it 113.51 & \it 167.91 & \it 542.60 & \it 969.60 \\
   11 & 31.21 &   0 & 3.06 & 9.80 & 22.12 & 19.25 & 25.90 & 29.95 & \it 44.51 & \it 43.62 & \it 52.70 & \it 80.13 & \it 65.11 & \it 77.25 & \it 362.10 & \it 674.10 \\
   12 & 35.07 &   0 & 2.28 & 6.03 & 21.28 & 21.10 & 30.41 & 29.77 & \it 46.67 & \it 44.84 & \it 57.17 & \it 78.81 & \it 66.58 & \it 80.18 & \it 339.80 & \it 681.10 \\
   13 & 25.97 &   0 & 0.00 & 6.18 & 6.36 & 22.09 & 17.67 & 21.55 & 17.31 & \it 42.93 & \it 47.00 & \it 66.43 & \it 61.66 & \it 74.91 & \it 335.00 & \it 579.50 \\
   14 & 17.58 &   0 & -2.20 & 4.81 & 7.46 & 17.41 & \it 21.72 & \it 23.22 & \it 25.21 & \it 36.65 & \it 61.53 & \it 72.31 & \it 57.71 & \it 101.33 & \it 283.40 & \it 576.60 \\
   15 & 48.59 &   0 & 3.17 & 32.75 & 30.81 & 41.20 & \it 50.88 & \it 66.37 & \it 72.36 & \it 107.39 & \it 137.15 & \it 111.80 & \it 155.11 & \it 194.37 & \it 753.70 & \it 1453.90 \\
   16 & 47.66 &   0 & 2.50 & 21.41 & 30.78 & 40.78 & \it 53.91 & \it 63.44 & \it 75.47 & \it 100.62 & \it 119.84 & \it 105.16 & \it 125.94 & \it 166.25 & \it 664.70 & \it 1288.40 \\
   17 & 22.88 &   0 & 4.03 & 9.93 & 16.69 & 18.85 & \it 29.35 & \it 29.06 & \it 43.45 & \it 53.96 & \it 60.86 & \it 62.01 & \it 57.55 & \it 88.92 & \it 512.80 & \it 675.10 \\
   18 & 23.06 &   0 & 3.25 & 11.88 & 13.01 & 16.83 & \it 25.60 & \it 28.15 & \it 50.50 & \it 49.65 & \it 51.91 & \it 66.20 & \it 66.34 & \it 71.43 & \it 462.00 & \it 669.70 \\
   19 & 36.22 &   0 & 4.02 & 20.95 & 10.94 & \it 39.01 & \it 41.28 & \it 41.18 & \it 52.94 & \it 72.86 & \it 96.80 & \it 73.37 & \it 96.39 & \it 145.30 & \it 643.00 & \it 1239.80 \\
   20 & 42.31 &   0 & 6.07 & 23.90 & 17.63 & 32.52 & \it 45.54 & \it 46.52 & 42.21 & \it 74.14 & \it 102.25 & \it 82.57 & \it 101.08 & \it 136.14 & \it 621.80 & \it 1202.90 \\
   21 & 78.18 &   0 & -6.23 & 10.10 & 28.34 & 53.09 & 56.35 & 66.12 & \it 82.41 & \it 120.85 & \it 150.16 & \it 134.85 & \it 173.29 & \it 254.72 & \it 886.30 & \it 2098.00 \\
   22 & 49.28 &   0 & 1.45 & 11.35 & 17.87 & 40.82 & 42.75 & \it 49.28 & \it 57.73 & \it 97.34 & \it 112.80 & \it 95.65 & \it 143.00 & \it 183.82 & \it 657.70 & \it 1516.70 \\
   23 & 67.29 &   0 & -3.23 & 11.11 & 29.58 & 60.09 & 53.83 & 65.57 & 52.90 & \it 113.15 & \it 164.95 & \it 104.85 & \it 186.07 & \it 260.09 & \it 897.70 & \it 2568.50 \\
   24 & 74.64 &   0 & 7.02 & 32.85 & 37.48 & \it 75.60 & \it 78.31 & \it 78.47 & 71.29 & \it 134.93 & \it 172.73 & \it 119.14 & \it 207.50 & \it 295.85 & \it 911.60 & \it 2689.30 \\
   25 & 53.64 &   0 & 1.82 & 3.44 & 25.30 & 37.05 & 34.41 & 47.77 & \it 90.28 & \it 100.00 & \it 121.86 & \it 114.37 & \it 112.75 & \it 140.69 & \it 725.90 & \it 1367.00 \\
   26 & 49.41 &   0 & 0.20 & 13.24 & 22.13 & 35.77 & 38.74 & 46.25 & \it 105.53 & \it 101.98 & \it 113.83 & \it 118.77 & \it 111.86 & \it 156.52 & \it 774.30 & \it 1335.60 \\
   27 & 62.45 &   0 & 13.36 & 14.06 & 37.41 & \it 72.86 & 60.34 & \it 77.64 & \it 120.39 & \it 141.77 & \it 159.77 & \it 130.24 & \it 175.25 & \it 240.79 & \it 1218.40 & \it 2251.30 \\
   28 & 87.63 &   0 & 15.72 & 23.14 & 54.59 & 83.84 & 83.99 & 75.84 & \it 154.73 & \it 135.23 & \it 165.50 & \it 179.18 & \it 184.43 & \it 308.15 & \it 1197.40 & \it 2291.40 \\
\midrule  
 \textbf{Rank} & 7 & 1 & 2 & 3 & 4 & 5 & 6 & 8 & 9 & 10 & 11 & 12 & 13 & 14 & 15 & 16  
\\ 
\bottomrule

\end{tabular}
}
\end{table}

\begin{tikzpicture}[overlay]
	\draw<2>[box1] (0.7,5.8) rectangle (1.4,0.55); 
	\draw<3>[box1] (6.1,5.8) rectangle (11.1,0.55);
	\draw<4>[box1] (0.1,0.6) rectangle (11.1,0.3)
	 node[title1] at (0.1,0.3) {LAHC classificado em 7\textsuperscript{\underline{o}}};      
	\end{tikzpicture}	

\end{frame}


%-----------------------------------------------------------------------------------------------
\begin{frame}{Resultados da etapa final da INRC-II}

\begin{columns}
\begin{column}{0.5\textwidth}
\begin{table}[!]
\centering
%\caption{Comparação do algoritmo proposto (LAHC) com os melhores resultados dos finalistas nas instâncias \textit{hidden}.Os resultados destacados em negrito são piores que a melhor solução encontrada pelo LAHC considerando 10 execuções.} 
\label{tab:hidden}
\definecolor{brown}{rgb}{0.28, 0.02, 0.03}
\renewcommand\it{\color{brown}\bf$^*$}
\resizebox{1\linewidth}{!}{
\begin{tabular}{rlrrrrrrrr}
  \toprule
id & LAHC & $A_1$ & $A_2$ & $A_3$ & $A_4$ & $A_5$ & $A_7$ & $A_6$ \\ 
  \midrule
   1 &  25.46 &   0 & 2.15 & 1.84 & 3.37 & 15.95 & 25.15 & 21.47 \\
    2 & 27.07 &   0 & -0.56 & 4.70 & 8.84 & 16.57 & 22.65 & 20.17 \\
    3 & 20.51 &   0 & 4.84 & 7.69 & 0.28 & \it 20.51 & 19.66 & 17.38 \\
    4 & 21.04 &   0 & -0.32 & 3.24 & 7.44 & \it 21.36 & \it 29.13 & 18.12 \\
    5 & 25.90 &   0 & -1.67 & 5.57 & 8.20 & 20.98 & 24.92 & \it 31.15 \\
    6 & 25.17 &   0 & -1.34 & 3.97 & 5.96 & 13.25 & 22.85 & \it 25.83 \\
    7 & 27.09 &   0 & 0.40 & 5.58 & 6.37 & 26.29 & \it 31.87 & 23.51 \\
    8 & 23.17 &   0 & 1.47 & 5.87 & 11.14 & 18.48 & \it 24.63 & \it 26.10 \\
    9 & 33.03 &   0 & 3.94 & 6.97 & 17.58 & 23.33 & \it 34.55 & 31.52 \\
   10 & 31.17 &   0 & 2.47 & 8.33 & 16.05 & 22.22 & 26.54 & 27.47 \\
   11 & - & \it   - & - & - & - & \it - & \it - & \it - \\
   12 &  33.10 &   0 & -4.69 & 15.69 & 17.41 & 31.72 & 29.14 & \it 41.03 \\
   13 &  37.11 &   0 & -3.42 & 16.03 & 11.15 & 31.88 & 34.15 & 36.06 \\
   14 &  34.52 &   0 & -0.18 & 19.57 & 13.17 & \it 38.08 & \it 34.70 & \it 39.86 \\
   15 &  42.61 &   0 & 1.06 & 14.44 & 21.13 & 31.34 & 38.56 & \it 44.37 \\
   16 &  33.44 &   0 & -4.81 & 7.05 & 12.95 & 29.02 & 26.89 & 31.15 \\
   17 &  34.46 &   0 & 2.53 & 16.39 & 19.59 & 34.12 & \it 38.51 & \it 43.92 \\
   18 &  43.69 &   0 & 3.91 & 23.45 & 20.25 & 36.77 & 41.03 & 40.85 \\
   19 &  32.51 &   0 & 1.64 & 16.75 & 10.02 & \it 33.66 & \it 35.47 & \it 34.15 \\
   20 &  35.43 &   0 & -5.52 & 18.50 & 15.01 & 31.76 & \it 35.95 & \it 39.44 \\
   21 &  28.10 &   0 & 3.88 & -0.19 & 14.23 & 23.48 & 27.73 & \it 34.94 \\
   22 &  24.28 &   0 & 3.70 & 3.50 & 14.81 & 21.40 & \it 29.01 & \it 31.07 \\
   23 &  25.05 &   0 & 4.65 & 2.42 & 15.76 & \it 26.67 & \it 32.73 & \it 37.58 \\
   24 &  30.80 &   0 & 6.37 & 7.39 & 17.25 & \it 32.44 & \it 38.40 & \it 47.64 \\
   25 &  29.53 &   0 & 6.25 & 6.25 & 17.46 & 26.08 & 29.09 & \it 39.44 \\
   26 &  22.96 &   0 & 3.70 & 4.26 & 12.41 & 20.56 & \it 25.93 & \it 26.67 \\
   27 &  26.59 &   0 & 6.55 & 10.32 & 14.29 & 23.41 & \it 32.54 & \it 39.68 \\
   28 &  25.43 &   0 & 0.38 & 2.49 & 16.63 & 21.61 & 23.33 & \it 40.54 \\
   29 &  25.79 &   0 & 3.74 & 5.91 & 13.19 & 25.00 & \it 27.95 & \it 34.84 \\
   30 &  28.49 &   0 & 2.49 & 5.16 & 13.77 & 19.31 & \it 31.17 & \it 30.21 \\
   \bottomrule
\end{tabular}
}
\end{table}

\end{column}
\begin{column}{0.5\textwidth}  %%<--- here
\begin{table}[!]
\centering
%\caption{Comparação do algoritmo proposto (LAHC) com os melhores resultados dos finalistas nas instâncias \textit{hidden}.Os resultados destacados em negrito são piores que a melhor solução encontrada pelo LAHC considerando 10 execuções.} 
\label{tab:hidden}
\definecolor{brown}{rgb}{0.28, 0.02, 0.03}
\renewcommand\it{\color{brown}\bf$^*$}
\resizebox{1\linewidth}{!}{
\begin{tabular}{rlrrrrrrrr}
  \toprule
id & LAHC & $A_1$ & $A_2$ & $A_3$ & $A_4$ & $A_5$ & $A_7$ & $A_6$ \\ 
  \midrule
   31 &  31.48 &   0 & 7.72 & 12.51 & 17.30 & \it 33.43 & \it 32.45 & 28.45 \\
   32 &  23.47 &   0 & 3.43 & 7.33 & 19.11 & \it 28.48 & \it 30.43 & 23.01 \\
   33 &  23.74 &   0 & 2.83 & 6.30 & 11.32 & \it 26.21 & 22.10 & 17.72 \\
   34 &  28.53 &   0 & 4.12 & 5.49 & 11.76 & 26.96 & \it 29.41 & 26.76 \\
   35 &  22.92 &   0 & 0.55 & 5.08 & 10.72 & \it 26.16 & \it 30.78 & 20.52 \\
   36 &  28.98 &   0 & 5.30 & 10.70 & 13.16 & 28.79 & \it 33.52 & 22.82 \\
   37 &  23.89 &   0 & 3.54 & 6.36 & 12.81 & \it 25.98 & \it 26.61 & 18.35 \\
   38 &  29.20 &   0 & 3.71 & 11.13 & 16.41 & 28.52 & \it 30.27 & 27.15 \\
   39 &  29.25 &   0 & 3.74 & 8.88 & 15.98 & 27.85 & \it 29.91 & 22.52 \\
   40 &  24.34 &   0 & 4.42 & 8.93 & 10.81 & \it 27.91 & \it 28.01 & 23.78 \\
   41 &  36.40 &   0 & -0.37 & 2.21 & 24.82 & 32.17 & 32.72 & \it 40.99 \\
   42 &  31.82 &   0 & -1.71 & 0.84 & 17.00 & 30.81 & \it 34.51 & \it 49.66 \\
   43 &  32.19 &   0 & -2.46 & -0.86 & 26.20 & \it 33.22 & \it 34.42 & \it 48.97 \\
   44 &  30.92 &   0 & -2.66 & 1.90 & 18.83 & \it 31.78 & \it 38.00 & \it 43.18 \\
   45 &  27.07 &   0 & -1.05 & 1.92 & 18.20 & \it 28.55 & \it 31.95 & \it 39.05 \\
   46 &  28.52 &   0 & 2.85 & 4.28 & 25.31 & \it 37.43 & \it 44.74 & \it 49.20 \\
   47 &  25.98 &   0 & -5.66 & 1.57 & 19.53 & \it 27.40 & \it 28.98 & \it 41.42 \\
   48 &  26.28 &   0 & -3.42 & 2.48 & 16.36 & \it 27.11 & \it 38.18 & \it 43.31 \\
   49 &  22.33 &   0 & -1.61 & 3.60 & 12.25 & \it 29.39 & \it 29.54 & \it 32.71 \\
   50 &  22.49 &   0 & -6.38 & 0.60 & 9.15 & 18.89 & \it 26.39 & \it 34.78 \\
   51 & 41.92 &   0 & -0.19 & 1.74 & 24.30 & 39.30 & 37.66 & \it 46.66 \\
   52 & 46.38 &   0 & -0.52 & 4.35 & 24.53 & 34.89 & 41.30 & \it 50.72 \\
   53 & 38.19 &   0 & -2.53 & 4.31 & 24.54 & 33.98 & 35.42 & \it 51.95 \\
   54 & 38.56 &   0 & -3.09 & -0.50 & 21.98 & 35.26 & \it 40.66 & \it 53.15 \\
   55 & 30.78 &   0 & -10.97 & -6.67 & 19.48 & 28.05 & \it 32.29 & \it 31.89 \\
   56 & 27.69 &   0 & -11.72 & -6.67 & 16.10 & 23.09 & 22.54 & \it 38.91 \\
   57 & 40.77 &   0 & -1.67 & -1.07 & 22.32 & 30.82 & 34.88 & \it 46.86 \\
   58 & 43.16 &   0 & -10.89 & -2.03 & 23.82 & 32.39 & 35.42 & \it 47.02 \\
   59 & 35.19 &   0 & -9.01 & 0.58 & 12.31 & 33.56 & \it 35.29 & \it 42.50 \\
   60 & 34.84 &   0 & -9.29 & -3.03 & 16.26 & 31.37 & \it 36.62 & \it 50.37 \\
 \midrule
  \textbf{Rank} & 6 & 1 & 2 & 3 & 4 & 5 & 7 & 8 \\
   \bottomrule
\end{tabular}
}
\end{table}

\end{column}
\end{columns}

\begin{tikzpicture}[overlay]
	\draw<2>[box1] (0.1,6.35) rectangle (0.8,0.3); 
    \draw<2>[box1] (6.35,6.35) rectangle (6.9,0.6); 
    \draw<3>[box1] (-0.25,4.2) rectangle (5,3.9); 
	\draw<4>[box1] (5.8,0.6) rectangle (11.1,0.3)
	 node[title1] at (5.8,0.3) {LAHC classificado em 6\textsuperscript{\underline{o}}};      
	\end{tikzpicture}	




\end{frame}

%-----------------------------------------------------------------------------------------------
\begin{frame}{Resultados comparando o LAHC com as melhores soluções conhecidas}

\begin{columns}
\begin{column}{0.5\textwidth}
\begin{table}[!]
\centering
%\caption{Comparação dos resultados obtidos pelos LAHC nas instâncias \textit{hidden} com os melhores valores conhecidos.} 
\label{tab:hiddenCV}
\resizebox{1\linewidth}{!}{
\begin{tabular}{rlrrrrrrrr}
  \toprule
id & Dimensão & $BKS$ & $\overline{obj}$ & $obj^*$ & \emph{gap} (\%) & $obj^-$ & $\sigma$ & $CV~{\scriptsize(\%)} $\\
  \midrule
1		&	140	&	1425	&	2133	&	2045	&	49.68	&	2260	&	70.29	&	3.30	\\
2		&	140	&	1615	&	2378.5	&	2300	&	47.28	&	2520	&	77.28	&	3.25	\\
3		&	140	&	1540	&	2190.5	&	2115	&	42.24	&	2320	&	72.36	&	3.30	\\
4		&	140	&	1365	&	2061.5	&	1870	&	51.03	&	2235	&	111.00	&	5.38	\\
5		&	140	&	1385	&	2081.5	&	1920	&	50.29	&	2300	&	94.58	&	4.54	\\
6		&	140	&	1335	&	2005.5	&	1890	&	50.22	&	2175	&	86.46	&	4.31	\\
7		&	140	&	1085	&	1737	&	1595	&	60.09	&	1805	&	63.39	&	3.65	\\
8		&	140	&	1525	&	2318.5	&	2100	&	52.03	&	2450	&	114.48	&	4.94	\\
9		&	140	&	1480	&	2345.5	&	2195	&	58.48	&	2485	&	107.97	&	4.60	\\
10		&	140	&	1620	&	2283.5	&	2125	&	40.96	&	2405	&	106.04	&	4.64	\\
11	&	280	&	3020	&	4215	&	3920	&	39.57	&	4400	&	170.00	&	4.03	\\
12	&	280	&	2770	&	4132	&	3860	&	49.17	&	4395	&	155.56	&	3.76	\\
13	&	280	&	2775	&	4144	&	3935	&	49.33	&	4315	&	112.09	&	2.70	\\
14	&	280	&	2805	&	3977	&	3780	&	41.78	&	4270	&	153.33	&	3.86	\\
15	&	280	&	2840	&	4261	&	4050	&	50.04	&	4515	&	158.75	&	3.73	\\
16	&	280	&	2910	&	4206	&	4070	&	44.54	&	4400	&	96.00	&	2.28	\\
17	&	280	&	2960	&	4235.5	&	3980	&	43.09	&	4400	&	130.24	&	3.07	\\
18	&	280	&	2815	&	4201	&	4045	&	49.24	&	4395	&	122.24	&	2.91	\\
19	&	280	&	3045	&	4220	&	4035	&	38.59	&	4500	&	152.15	&	3.61	\\
20	&	280	&	2715	&	4061	&	3880	&	49.58	&	4325	&	119.18	&	2.93	\\
21		&	280	&	2460	&	3623.5	&	3465	&	47.30	&	3765	&	76.60	&	2.11	\\
22		&	280	&	2330	&	3216	&	3020	&	38.03	&	3335	&	92.73	&	2.88	\\
23		&	280	&	2315	&	3232	&	3095	&	39.61	&	3370	&	98.66	&	3.05	\\
24		&	280	&	2400	&	3244.5	&	3185	&	35.19	&	3315	&	40.38	&	1.24	\\
25		&	280	&	2225	&	3218	&	3005	&	44.63	&	3385	&	114.94	&	3.57	\\
26		&	280	&	2615	&	3397	&	3320	&	29.90	&	3535	&	64.17	&	1.89	\\
27		&	280	&	2415	&	3315	&	3190	&	37.27	&	3615	&	123.17	&	3.72	\\
28		&	280	&	2405	&	3370	&	3280	&	40.12	&	3460	&	66.75	&	1.98	\\
29		&	280	&	2390	&	3276	&	3195	&	37.07	&	3355	&	47.60	&	1.45	\\
30		&	280	&	2480	&	3454.5	&	3360	&	39.29	&	3550	&	64.53	&	1.87	\\
  \midrule		
  \multicolumn{4}{l}{Média das instâncias pequenas}   -       &   -       &	44.85   &   -       &   -       &   3.29  \\		
   \bottomrule
\end{tabular}
}
\end{table}

\end{column}
\begin{column}{0.5\textwidth}  %%<--- here
\begin{table}[!]
\centering
%\caption{Comparação dos resultados obtidos pelos LAHC nas instâncias \textit{hidden} com os melhores valores conhecidos.} 
\label{tab:hiddenCV}
\resizebox{1\linewidth}{!}{
\begin{tabular}{rlrrrrrrrr}
  \toprule
id & Dimensão & $BKS$ & $\overline{obj}$ & $obj^*$ & \emph{gap} (\%) & $obj^-$ & $\sigma$ & $CV~{\scriptsize(\%)} $\\
  \midrule
41	&	440	&	2560	&	3865	&	3710	&	50.98	&	4135	&	141.55	&	3.66	\\
42	&	440	&	2640	&	4162	&	3915	&	57.65	&	4345	&	139.82	&	3.36	\\
43	&	440	&	2690	&	4035	&	3860	&	50.00	&	4180	&	106.56	&	2.64	\\
44	&	440	&	2705	&	3946	&	3790	&	45.88	&	4175	&	113.10	&	2.87	\\
45	&	440	&	3170	&	4405.5	&	4295	&	38.97	&	4500	&	69.70	&	1.58	\\
46	&	440	&	2630	&	3948	&	3605	&	50.11	&	4160	&	149.13	&	3.78	\\
47	&	440	&	2960	&	4189.5	&	4000	&	41.54	&	4300	&	93.61	&	2.23	\\
48	&	440	&	2770	&	3998.5	&	3820	&	44.35	&	4125	&	83.70	&	2.09	\\
49	&	440	&	3140	&	4394	&	4245	&	39.94	&	4460	&	70.15	&	1.60	\\
50	&	440	&	3005	&	4222.5	&	4085	&	40.52	&	4350	&	83.57	&	1.98	\\
31&	560	&	5115	&	6847	&	6725	&	33.86	&	7195	&	140.85	&	2.06	\\
32&	560	&	5390	&	7064	&	6655	&	31.06	&	7310	&	200.44	&	2.84	\\
33&	560	&	5475	&	7092.5	&	6775	&	29.54	&	7350	&	206.41	&	2.91	\\
34&	560	&	5100	&	6773	&	6555	&	32.80	&	7025	&	141.05	&	2.08	\\
35&	560	&	5410	&	6946	&	6650	&	28.39	&	7230	&	198.28	&	2.85	\\
36&	560	&	5280	&	7012.5	&	6810	&	32.81	&	7370	&	162.74	&	2.32	\\
37&	560	&	5505	&	7061.5	&	6820	&	28.27	&	7295	&	173.54	&	2.46	\\
38&	560	&	5120	&	6855	&	6615	&	33.89	&	7070	&	157.25	&	2.29	\\
39&	560	&	5350	&	7180	&	6915	&	34.21	&	7445	&	167.84	&	2.34	\\
40&	560	&	5320	&	6780	&	6615	&	27.44	&	7015	&	120.71	&	1.78	\\
 \midrule		
 \multicolumn{4}{l}{Média das instâncias médias}   -       &   -       &	38.61   &   -       &   -       &   2.49  \\		
  \midrule	
51	&	880	&	5155	&	7412	&	7330	&	43.78	&	7520	&	68.36	&	0.92	\\
52	&	880	&	4805	&	7355.5	&	7070	&	53.08	&	7575	&	163.15	&	2.22	\\
53	&	880	&	4750	&	7045.5	&	6730	&	48.33	&	7410	&	188.98	&	2.68	\\
54	&	880	&	4855	&	7197.5	&	6935	&	48.25	&	7455	&	173.64	&	2.41	\\
55	&	880	&	4465	&	6919	&	6480	&	54.96	&	7180	&	201.90	&	2.92	\\
56	&	880	&	4865	&	7321	&	6940	&	50.48	&	7615	&	179.56	&	2.45	\\
57	&	880	&	5090	&	7590	&	7285	&	49.12	&	7805	&	151.58	&	2.00	\\
58	&	880	&	4315	&	6988	&	6850	&	61.95	&	7110	&	92.56	&	1.32	\\
59	&	880	&	4770	&	7256.5	&	7030	&	52.13	&	7520	&	155.77	&	2.15	\\
60	&	880	&	4360	&	6778	&	6425	&	55.46	&	7080	&	189.85	&	2.80	\\
\midrule		
\multicolumn{4}{l}{Média das instâncias grandes}  &   -  &	51.75   &   -       &   -       &   2.18  \\		
	 \midrule								
 \multicolumn{4}{l}{ \textbf{Média Geral} }          &  -        & 45.07     &   -       &  -        &    2.82 \\ 
   \bottomrule
\end{tabular}
}
\end{table}
    
\end{column}
\end{columns}
\begin{tikzpicture}[overlay]
	\draw<2>[box1] (0.1,6.25) rectangle (0.9,0.8); 
	\draw<2>[box1] (6.15,6.5) rectangle (7.,0.8) 
	 node[title1] at (0.1,0.8) {Dimensão = n\textsuperscript{\underline{o}} enfermeiros x n\textsuperscript{\underline{o}} semanas};   
	\draw<3>[box1] (-0.3,0.55) rectangle (5.1,0.9);
    \draw<3>[box1] (5.7,2.8) rectangle (11.1,2.55);
    \draw<3>[box1] (5.7,0.55) rectangle (11.1,0.85);
    \draw<4>[box1] (5.7,0.3) rectangle (11.1,0.6);
	\end{tikzpicture}	
\end{frame}

%-------------------------------------------------------------------------------------------------
\section{Considerações Finais}

\begin{frame}{Conclusões}

\begin{enumerate}
	\item \textbf{Principais contribuições:} modelo matemático e o algoritmo baseado no LAHC;
		\item Modelo proposto possui uma \structure{descrição mais fácil} de compreender que o apresentado na definição da INRC-II;
		\item Algoritmo é \structure{robusto} em relação à variação dos resultados e produziu boas soluções para instâncias de \structure{dimensões médias} com uma implementação relativamente simples;
		\item No ajuste de parâmetros do LAHC, dos \structure{sete movimentos propostos}, cinco apresentaram uma utilidade significativa;
		\item Na primeira etapa da INRC-II, o LAHC \structure{se classificou na sétima colocação};
		\item Na etapa final, o LAHC \structure{se classificou em sexto lugar}, de um total de quinze competidores internacionais, o que consideramos um resultado competitivo.
\end{enumerate}
\end{frame}

%--------------------------------------------------------------------------------------------------

\begin{frame}{Trabalhos Futuros}
	\begin{enumerate}
		\item Ajuste dinâmico dos parâmetros conforme dimensão da instância;
		\item Avaliação de novos movimentos e inclusão de uma Busca Local no final do LAHC;
		\item Linearização do modelo matemático e avaliação dos métodos exatos;
		\item Realização de estudo de caso em um ambiente hospitalar real com o objetivo de coletar instâncias reais e analisar qualitativamente a adequação do método proposto.
	\end{enumerate}

\end{frame}

%--------------------------------------------------------------------------------------------------

\def\newblock{\hskip .9em plus .31em minus .05em}
\bibliographystyle{sbc}     

\begin{frame}[allowframebreaks]{Referências}
\bibliography{refs}   
\end{frame}


%-------------------------------------------------------------------------------------------------
\begin{frame}[noframenumbering]
   \begin{center}
   	
   	 \medskip
   
      {\Huge Obrigada!}\\
      ~
      
      \textbf{Victória S. Portella}
      ~

      vsportella@inf.ufrgs.br\\
   
   \medskip
      
   \end{center}
    \begin{center}   	
    	
	\end{center}
  
\end{frame}

\section*{Slides complementares}
%----------------------------------------------------------------------------------------------
%\begin{frame}[noframenumbering]{Restrição com arcos}
	
%\end{frame}


\end{document}




